\documentclass{article}

\usepackage{tabularx}
\usepackage{booktabs}

\title{Problem Statement and Goals\\\progname}

\author{\authname}

\date{}

\input{../Common.text}
\renewcommand{\progname}{LingoLyrics}
\renewcommand{\authname}{Team LingoLyrics
\\ Rhea Gokhale
\\ Omar Bakr
\\ Raymond Ma
\\ Aansh Kotian}

\begin{document}

\maketitle

\begin{table}[hp]
\caption{Revision History} \label{TblRevisionHistory}
\begin{tabularx}{\textwidth}{llX}
\toprule
\textbf{Date} & \textbf{Developer(s)} & \textbf{Change}\\
\midrule
September 22, 2026 & Omar Bakr & Initial Revision 0 draft\\
September 28, 2026 & Raymond Ma & Filled out the template; implemented supervisor feedback; reworded documentation with Codex.\\
September 28, 2026 & Rhea Gokhale & Changed goals to selling point level, confirmed customer documents, overall polish and review \\
\bottomrule
\end{tabularx}
\end{table}

\section{Problem Statement}
This section describes the learning and content-authoring problem addressed by
LingoLyrics, its high-level inputs and outputs, intended stakeholders, and
operating environment.

\subsection{Problem}
Self-directed language learners at an early-intermediate to intermediate level
often use songs to encounter authentic vocabulary, idioms, and cultural
context. However, conventional music players and lyric websites do not help
learners turn those moments into focused, in-context study. When a learner
encounters an unfamiliar word or phrase, they must pause, search elsewhere for
meaning, and manually transfer useful vocabulary into a study system. This
interrupts the learning experience and separates vocabulary from the lyrical
context needed to understand its intended meaning.

The challenge is greater for languages with non-Latin scripts, different word
boundaries, or different sentence structures. Learners may require
language-appropriate segmentation, transliteration, and contextual explanations
in addition to direct translation. Creators also lack a consistent, reusable way
to prepare synchronized lyric lessons with timing, translations, and
annotations.

LingoLyrics will deliver a mobile song-based learning experience and a companion
content-authoring workflow. Together, they will allow learners to study
synchronized, annotated lyrics and retain selected vocabulary through Anki.

\subsection{Key Terms}
\begin{itemize}
  \item \textbf{Song map:} a structured lesson file that connects an audio
        track with timed lyrics, translations, transliterations, annotations,
        and language metadata.
  \item \textbf{Contextual annotation:} an explanation of a word, phrase, or
        lyric line that captures its intended meaning in the song.
  \item \textbf{Segmentation:} splitting a line of text into individual words,
        which is needed for languages such as Japanese that do not put spaces
        between words.
  \item \textbf{Transliteration:} writing a word in another alphabet (for
        example, Japanese written in Latin letters) so a learner can
        pronounce it before they can read the original script.    
  \item \textbf{Anki:} a spaced-repetition flashcard application used for
        vocabulary review.
\end{itemize}

\subsection{Inputs and Outputs}
\textbf{Inputs:}
\begin{itemize}
  \item Audio tracks, lyric text, timed word and line boundaries, translations,
        transliterations, and contextual annotations.
  \item Learner interactions, display preferences, and selected vocabulary.
\end{itemize}

\textbf{Outputs:}
\begin{itemize}
  \item A validated, reusable song map containing synchronized lyric and
        annotation layers.
  \item A synchronized lesson view with contextual support for selected words
        or lines.
  \item An optional vocabulary item, including its selected context, sent to
        the learner's Anki workflow.
\end{itemize}

\subsection{Stakeholders}
The primary stakeholders are self-directed early-intermediate to intermediate
language learners, who need an engaging way to study authentic language in
context, and song-map creators, who need a practical workflow for accurately
synchronizing and annotating lessons.

Secondary stakeholders include:
\begin{itemize}
  \item Anki users, who benefit from reliable vocabulary handoff;
  \item the project supervisor and development team;
  \item copyright holders and content providers, whose rights constrain what
        material may be used or distributed.
\end{itemize}

\subsection{Environment}
The learner-facing application is designed for short, self-directed study
moments on an iOS or Android smartphone, including while listening to music
with headphones. The creator workflow will be delivered through a web-based
editor used on a laptop or desktop computer. Both workflows must use a shared
song-map representation.

The initial release will use team-created, public-domain, or appropriately
licensed content. It does not promise streaming-service integration, a
production music catalogue, universal language support, or distribution rights
for copyrighted lyrics or recordings.

\section{Goals}
The table below contains the goals that LingoLyrics must achieve to be
considered acceptable and complete.

Language learners often have to reconstruct a song's context across separate
tools, such as lyrics websites, translation tools, and music players. This
interrupts listening and makes it harder to understand and retain vocabulary
in its original lyrical context.

\progname{} combines timed lyrics, translations, contextual annotations, and
vocabulary export in one workflow. Its reusable song-map format prevents
learners and teachers from having to reconstruct the same context separately.

Table~\ref{TblGoals} lists the goals that \progname{} must achieve to be
considered acceptable and complete. Each goal has a measure that will be
refined into requirements in the Software Requirements Specification (SRS).

\begin{table}[hp]
\caption{Minimum Viable Product Goals}
\begin{tabularx}{\textwidth}{lXX} \\[-8pt]
\toprule
\textbf{Goal} & \textbf{Explanation} & \textbf{Reasoning}\\
\midrule
Synchronized Lesson Playback & Across every lyric-line transition in a fixed two-song acceptance suite, at least 95\% occur within 0.5 seconds of their mapped timestamps. & Timing keeps the lesson connected to the music.\\
In-Context Vocabulary Support & Every annotated target in the fixed two-song acceptance suite displays its translation and contextual annotation when selected. & Meaning is clearer when vocabulary is preserved in context.\\
Creator-to-Learner Interoperability & A song map created by the editor validates against one shared format and is opened by the learner application without manual file modification. & Learners receive consistent lessons, while creators avoid duplicate content preparation.\\
Practical Lesson Creation & At least two of three test creators can produce a complete map for a pre-defined short test track within 45 minutes, without manually editing the underlying map file; each resulting map must validate and play in the learner application. & A lesson system is only useful if content can be prepared consistently.\\
Language-Aware Support & Demonstrate two languages with different text-handling needs; at least one must require non-Latin rendering, language-specific segmentation, or transliteration. & This validates that the product is not limited to one writing system.\\
Vocabulary Retention Handoff & Sending a selected vocabulary item creates a note in the named \emph{LingoLyrics Test} Anki deck with correctly populated \emph{Term}, \emph{Translation}, and \emph{Context} fields. & This connects song learning to spaced-repetition study.\\
\bottomrule
\end{tabularx}
\end{table}

For goal evaluation, the \textbf{fixed two-song acceptance suite} contains one
map for each demonstrated language and is defined before implementation testing
begins. An \textbf{annotated target} is a word or line with a translation and
contextual annotation. A \textbf{complete map} contains all required audio,
timing, lyric, and annotation data. A \textbf{test creator} is a participant
given the same supplied audio, lyric, and annotation material; the
\textbf{pre-defined short test track} contains no more than 12 lyric lines.

\newpage

\section{Stretch Goals}
The table below shows goals that are not required for the product to be viable,
but would provide additional value if achieved.

\begin{table}[hp]
\caption{Stretch Goals}
\begin{tabularx}{\textwidth}{lXX} \\[-8pt]
\toprule
\textbf{Goal} & \textbf{Explanation} & \textbf{Reasoning}\\
\midrule
Third-Language Support & Add another language and a pluggable language-processing approach. & Demonstrates that the format can scale.\\
Song-Map Import and Export & Add migration and interchange tooling for song maps. & Helps maps remain reusable over time.\\
Collaborative Review & Support sharing or review of song maps. & Encourages higher-quality learning content.\\
Learner Progress and Accessibility & Add progress tracking, in-app review, or accessibility enhancements. & Extends learning value beyond playback.\\
Timing and Annotation Assistance & Evaluate automatic assistance while preserving creator verification. & May reduce authoring effort without sacrificing accuracy.\\
\bottomrule
\end{tabularx}
\end{table}

\newpage

\section{Custom Documents (CDs)}
\begin{itemize}
  \item Fall: Design Thinking Report or User/Stakeholder/Client Interview Report: interviews
        with target learners and song-map creators to validate the problem and
        the learner and creator workflows before design begins.
  \item Winter: User Manual or Usability Report: usability testing of the mobile
        learner app and the web editor with representative users, measured
        against the goals in Table~\ref{TblGoals}.
\end{itemize}

The Winter CD was changed from a User Manual to a Usability Report following
instructor feedback. The fall custom document will validate learner and creator workflows. The
winter custom document will document and evaluate the editor and mobile lesson
flows.

\section*{Generative AI Use}
The Generative AI Use Disclosure covering this document is in the appendix of
the \progname{} Development Plan.

\newpage{}

\section*{Appendix --- Reflection}
\begin{enumerate}
  \item \textbf{What went well while writing this deliverable?}\\
  The team quickly aligned on the project direction and maintained open,
  consistent communication while preparing the document. This made it easier to
  discuss decisions, divide work, and keep the problem statement and goals
  consistent.

  \item \textbf{What pain points did you experience during this deliverable,
  and how did you resolve them?}\\
  A few meetings had to be rescheduled, which made coordination harder; we
  resolved this by setting recurring Friday meeting times. The harder problem
  was the level of detail in the goals. Our first drafts were either too vague
  to measure or so specific that they read like test cases (for example, naming
  the exact Anki deck and field names). We resolved this by checking each goal
  against the checklist question ``is this a selling point or a requirement?'',
  keeping the measurable part, and moving implementation detail to the SRS. We
  also found that parts of our drafts assumed a web app while others assumed a
  phone app; we agreed on a mobile learner app with a web-based editor and
  updated both documents to match.

  \item \textbf{How did you and your team adjust the scope of your goals to
  ensure they are suitable for a Capstone project?}\\
  Professor feedback identified ``language learners'' as too broad a target
  group, so we narrowed the intended users to self-directed, early
  to intermediate learners who use songs to study a target language. We then
  limited the core product to two languages, a two-song acceptance suite, and
  team-created or licensed content, and explicitly excluded streaming-service
  integration and a production music catalogue. Features we liked but could not
  guarantee, such as a third language, collaborative review, and automatic
  timing assistance, became stretch goals. Keeping one non-Latin language in
  the core goals keeps the project technically challenging enough for a senior
  design project.
\end{enumerate}

\end{document}
